% ============================================================================
%  Formulario T-15. Nivelación de recursos. Archivo propio
%  AUDIT-Formulario-T-15.pdf (Comunicado 10, sección 1).
%  Escrito por la dupla 4 por encargo del usuario (D4-61, D4-64 y D4-67):
%  dotación de los meses 13 a 15, asignación por rol, estimaciones de tres
%  puntos y malla CPM. Las filas de las tablas salen del mismo cálculo que las
%  figuras del S7 (D4/lucid/d11_pert.py). Las horas hombre y el resto de las
%  etapas son de la dupla 2 y llevan su marcador.
% ============================================================================
\begin{formulario}{T-15}

La \tabref{t15} resume cada etapa del \art{17.1}{12}. Las filas de la marcha blanca de la Etapa 1 y
del desarrollo de la Etapa 2 traen la dotación de los meses 13 a 15, que el \art{17.2}{13} pide
demostrar en este formulario.

\begin{formularioTQuince}
\etapaTQuince{etapa = {Etapa 1, desarrollo}, hh = {\requiereDupla{2}{HH}}, personas = {\requiereDupla{2}{peak}}, frentes = {\requiereDupla{2}{frentes}}, meses = {1 a 12}}
\etapaTQuince{etapa = {Etapa 1, marcha blanca}, hh = {\requiereDupla{2}{HH}}, personas = {26}, frentes = {6}, meses = {13 a 15}}
\etapaTQuince{etapa = {Etapa 2, desarrollo}, hh = {\requiereDupla{2}{HH}}, personas = {18 en los meses 13 a 15}, frentes = {2}, meses = {13 a 18}}
\etapaTQuince{etapa = {Etapa 2, marcha blanca}, hh = {\requiereDupla{2}{HH}}, personas = {\requiereDupla{2}{peak}}, frentes = {\requiereDupla{2}{frentes}}, meses = {19 y 20}}
\etapaTQuince{etapa = {Operación}, hh = {\requiereDupla{2}{HH}}, personas = {\requiereDupla{2}{peak}}, frentes = {\requiereDupla{2}{frentes}}, meses = {21 a 56}}
\end{formularioTQuince}

En los meses 13 a 15 trabajan 37 personas. Las 26 de la Escuadra de Estabilización E1 y las 18 de la
Escuadra de Construcción E2 suman 44 porque los siete roles de dirección están en las dos. E1 trabaja
en seis frentes, los cinco terminales y la guardia central, y E2 en dos, software y firmware.
\requiereDupla{2}{horas hombre por tarea, paquete y etapa, curvas de horas hombre por etapa y de la
continuidad operacional, y peak de personas de los meses 16 a 18}

La \tabref{7-asignacion} reparte la capacidad de cada rol entre las dos escuadras en los meses 13 a 15,
según la regla de asignación de la \verSeccion{sec:7-particion}. Cada valor está en FTE, una persona a
jornada completa, y entre paréntesis va el porcentaje sobre la capacidad del rol. «Planta» es personal
de los 22 profesionales del Subdocumento~1 y «refuerzo» es personal que audIT incorpora para el
proyecto.

\begin{tabla}[ancho=completo, fuente=condensada]{Asignación de capacidad por rol en los meses 13 a 15}{7-asignacion}{X R{22mm} R{30mm} R{30mm} R{24mm}}
Rol (personas, origen) & Capacidad (FTE) & Escuadra E1 & Escuadra E2 & Total asignado \\ \midrule
Jefe de Proyecto (1, planta) & 1,0 & 0,4 (40~\%) & 0,4 (40~\%) & 80~\% \\
Arquitecto de Solución (1, planta) & 1,0 & 0,3 (30~\%) & 0,5 (50~\%) & 80~\% \\
Encargado de Seguridad de la Información (1, planta) & 1,0 & 0,4 (40~\%) & 0,4 (40~\%) & 80~\% \\
Líder de Calidad y Pruebas (1, planta) & 1,0 & 0,3 (30~\%) & 0,5 (50~\%) & 80~\% \\
Líder de Datos (1, planta) & 1,0 & 0,3 (30~\%) & 0,5 (50~\%) & 80~\% \\
Líder de Operación / SRE (1, planta) & 1,0 & 1,0 (100~\%) & 0,0 (0~\%) & 100~\% \\
Analista PMO (3, planta) & 3,0 & 2,0 (67~\%) & 1,0 (33~\%) & 100~\% \\
Líder de Desarrollo (1, planta) & 1,0 & 0,0 (0~\%) & 1,0 (100~\%) & 100~\% \\
Ingeniero de software sénior (2, planta) & 2,0 & 1,0 (50~\%) & 1,0 (50~\%) & 100~\% \\
Ingeniero frontend y móvil (2, planta) & 2,0 & 1,0 (50~\%) & 1,0 (50~\%) & 100~\% \\
Ingeniero de datos (1, planta) & 1,0 & 1,0 (100~\%) & 0,0 (0~\%) & 100~\% \\
Ingeniero de pruebas y automatización (1, planta) & 1,0 & 1,0 (100~\%) & 0,0 (0~\%) & 100~\% \\
Ingeniero de firmware (3, planta) & 3,0 & 1,0 (33~\%) & 2,0 (67~\%) & 100~\% \\
Ingeniero electrónico (1, planta) & 1,0 & 0,0 (0~\%) & 1,0 (100~\%) & 100~\% \\
Ingeniero de telecomunicaciones y campo (1, planta) & 1,0 & 1,0 (100~\%) & 0,0 (0~\%) & 100~\% \\
Líder Funcional (1, refuerzo) & 1,0 & 0,4 (40~\%) & 0,4 (40~\%) & 80~\% \\
Líder de Integración (1, refuerzo) & 1,0 & 0,4 (40~\%) & 0,4 (40~\%) & 80~\% \\
Líder de Implantación y Gestión del Cambio (1, refuerzo) & 1,0 & 1,0 (100~\%) & 0,0 (0~\%) & 100~\% \\
Analista de implantación en terreno (7, refuerzo) & 7,0 & 7,0 (100~\%) & 0,0 (0~\%) & 100~\% \\
Ingeniero SRE / DevOps (2, refuerzo) & 2,0 & 2,0 (100~\%) & 0,0 (0~\%) & 100~\% \\
Desarrollador full-stack (2, refuerzo) & 2,0 & 0,0 (0~\%) & 2,0 (100~\%) & 100~\% \\
Ingeniero de datos y analítica (1, refuerzo) & 1,0 & 0,0 (0~\%) & 1,0 (100~\%) & 100~\% \\
Analista de QA (1, refuerzo) & 1,0 & 0,0 (0~\%) & 1,0 (100~\%) & 100~\% \\
Total (37 personas: 21 de planta y 16 de refuerzo) & 37,0 & 21,5 (58~\%) & 14,1 (38~\%) & 96~\% \\
\end{tabla}
\notaTabla{Fuente: elaboración propia.}

Ninguna fila supera el 100 \% de la capacidad del rol, así que no hay sobreasignación. Las 16 filas de
roles dedicados suman exactamente su capacidad. Las siete filas de roles de dirección suman 80 \% y
dejan 1,4 FTE para coordinar entre escuadras, por eso el total asignado es 96 \% y no 100 \%. Las
cuentas cuadran: 21,5 FTE en E1, 14,1 en E2 y 1,4 de coordinación suman los 37,0 FTE del equipo. En los
tres meses eso equivale a 64,5 meses-persona en E1, 42,3 en E2 y 4,2 de coordinación, 111 en total.

La \verSeccion{sec:7-ruta} calcula la ruta crítica y las holguras de la implementación con una red de
25 actividades en días hábiles. La \tabref{7-tres-puntos} trae las estimaciones de tres puntos de 18
de ellas: optimista (a), más probable (m) y pesimista (b), con la duración esperada y la varianza. Las
otras siete son las dos reservas, las dos marchas blancas, los dos pasos a producción y la
estabilización de la Etapa 1, cuya duración fija el contrato o declara audIT.

\begin{tabla}[ancho=completo, fuente=condensada]{Estimaciones de tres puntos, en días hábiles}{7-tres-puntos}{L{14mm} R{18mm} R{24mm} R{18mm} R{30mm} X}
Actividad & a & m & b & Duración esperada & Varianza \\ \midrule
A01 & 4 & 5 & 6 & 5 & 0,11 \\
A02 & 28 & 32 & 42 & 33 & 5,44 \\
A03 & 38 & 44 & 56 & 45 & 9,00 \\
A04 & 45 & 58 & 83 & 60 & 40,11 \\
A05 & 32 & 37 & 48 & 38 & 7,11 \\
A06 & 95 & 118 & 153 & 120 & 93,44 \\
A07 & 29 & 35 & 47 & 36 & 9,00 \\
A08 & 50 & 59 & 74 & 60 & 16,00 \\
A09 & 70 & 78 & 98 & 80 & 21,78 \\
A10 & 56 & 65 & 80 & 66 & 16,00 \\
A11 & 50 & 59 & 74 & 60 & 16,00 \\
A12 & 30 & 33 & 42 & 34 & 4,00 \\
A13 & 36 & 40 & 44 & 40 & 1,78 \\
A16 & 42 & 50 & 58 & 50 & 7,11 \\
A17 & 31 & 35 & 45 & 36 & 5,44 \\
A18 & 100 & 118 & 148 & 120 & 64,00 \\
A20 & 42 & 48 & 66 & 50 & 16,00 \\
A22 & 20 & 23 & 32 & 24 & 4,00 \\
\end{tabla}
\notaTabla{Fuente: elaboración propia. Duración esperada (a + 4m + b) / 6 y varianza ((b − a) / 6)².}

Las varianzas más altas son las de la adhesión de los transportistas (A06), la instalación progresiva a
bordo (A18) y la factibilidad con los proveedores (A04). Las tres dependen de terceros y ninguna está
en la ruta crítica.

La \tabref{7-malla} es la malla de precedencias, con el inicio y el término tempranos (ES y EF), el
inicio y el término tardíos (LS y LF), la holgura total (HT) y la holgura libre (HL) de cada actividad.
La duración va en días hábiles y el día 0 es el inicio del mes 1.

\begin{tabla}[ancho=completo, fuente=condensada]{Malla de precedencias y cálculo de holguras}{7-malla}{L{8mm} X L{17mm} R{11mm} R{7mm} R{7mm} R{7mm} R{7mm} R{6mm} R{6mm} C{11mm}}
ID & Descripción & Predecesoras & Duración & ES & EF & LS & LF & HT & HL & Ruta crítica \\ \midrule
A01 & Inicio del proyecto: plan de dirección aprobado y equipo constituido & Ninguna & 5 & 0 & 5 & 0 & 5 & 0 & 0 & Sí \\
A02 & Levantamiento de procesos en los cinco terminales, línea base de alcance y matriz de trazabilidad (H1) & A01 & 33 & 5 & 38 & 5 & 38 & 0 & 0 & Sí \\
A03 & Caracterización en terreno de la cobertura móvil de las rutas & A01 & 45 & 5 & 50 & 87 & 132 & 82 & 62 & No \\
A04 & Verificación de factibilidad con los tres proveedores de posicionamiento y con los fabricantes de los camiones con telemetría de fábrica & A01 & 60 & 5 & 65 & 66 & 126 & 61 & 47 & No \\
A05 & Arquitectura, plan de seguridad y modelo de datos (H2) & A02 & 38 & 38 & 76 & 38 & 76 & 0 & 0 & Sí \\
A06 & Plan de adhesión de los 148 transportistas subcontratados: contrato, incentivo y piloto & A02 & 120 & 38 & 158 & 100 & 220 & 62 & 34 & No \\
A07 & Infraestructura híbrida y ambientes DEV, QA, PREPROD y PROD con observabilidad (H3) & A05 & 36 & 76 & 112 & 76 & 112 & 0 & 0 & Sí \\
A08 & Migración y verificación documental de las 6.000 vigencias de cuatro planillas & A05, A07 & 60 & 112 & 172 & 132 & 192 & 20 & 20 & No \\
A09 & Construcción de la Etapa 1 en ocho sprints de dos semanas y entrega para pruebas (H4) & A07 & 80 & 112 & 192 & 112 & 192 & 0 & 0 & Sí \\
A10 & Integraciones con el sistema de gestión de transporte de 2013, el sistema contable y las tres plataformas de posicionamiento & A04, A07 & 66 & 112 & 178 & 126 & 192 & 14 & 14 & No \\
A11 & Piloto de equipamiento a bordo por familia de camión, instalado en terminal camión por camión & A03, A07 & 60 & 112 & 172 & 132 & 192 & 20 & 20 & No \\
A12 & Pruebas integrales y certificación de la Etapa 1: aceptación de usuario, carga, resiliencia y seguridad ofensiva (H5) & A08, A09, A10, A11 & 34 & 192 & 226 & 192 & 226 & 0 & 0 & Sí \\
A13 & Capacitación de conductores en los terminales y de la torre 24x7 & A06, A09 & 40 & 192 & 232 & 220 & 260 & 28 & 8 & No \\
A14 & Reserva de contingencia de cronograma de la Etapa 1 & A12 & 14 & 226 & 240 & 226 & 240 & 0 & 0 & Sí \\
A15 & Marcha blanca de la Etapa 1, meses 13 a 15, con medición diaria y reversión activa (H6 al inicio) & A13, A14 & 60 & 240 & 300 & 260 & 320 & 20 & 0 & No \\
A16 & Transferencia tecnológica y certificación del equipo de tecnologías de información del mandante & A14 & 50 & 240 & 290 & 270 & 320 & 30 & 10 & No \\
A17 & Levantamiento y diseño detallado de la Etapa 2 (H8) & A14 & 36 & 240 & 276 & 240 & 276 & 0 & 0 & Sí \\
A18 & Escalamiento telemático e instalación progresiva a bordo & A06, A11, A14 & 120 & 240 & 360 & 280 & 400 & 40 & 40 & No \\
A19 & Paso a producción de la Etapa 1 y acta de cierre de la marcha blanca (H7) & A15, A16 & 10 & 300 & 310 & 320 & 330 & 20 & 0 & No \\
A20 & Construcción de la Etapa 2 en cinco sprints de dos semanas y entrega para pruebas (H9) & A17 & 50 & 276 & 326 & 276 & 326 & 0 & 0 & Sí \\
A21 & Estabilización posterior al paso a producción de la Etapa 1, con presencia en terminales en el relevo & A19 & 30 & 310 & 340 & 330 & 360 & 20 & 20 & No \\
A22 & Certificación de la Etapa 2 y cierre del desarrollo (H10) & A20 & 24 & 326 & 350 & 326 & 350 & 0 & 0 & Sí \\
A23 & Reserva de contingencia de cronograma de la Etapa 2 & A22 & 10 & 350 & 360 & 350 & 360 & 0 & 0 & Sí \\
A24 & Marcha blanca de la Etapa 2 en convivencia con la Etapa 1 en producción (H11 al inicio) & A21, A23 & 40 & 360 & 400 & 360 & 400 & 0 & 0 & Sí \\
A25 & Paso a producción de la Etapa 2 y aceptación final del proyecto (H12) & A18, A24 & 20 & 400 & 420 & 400 & 420 & 0 & 0 & Sí \\
\end{tabla}
\notaTabla{Fuente: elaboración propia.}

Las 13 actividades de la ruta crítica tienen holgura total y libre igual a cero, y la red termina el día
hábil 420, fin del mes 21, que es el hito H12. La \verSeccion{sec:7-ruta} analiza la ruta y las holguras.

\end{formulario}
